\documentclass[11pt]{article}
\usepackage{amsmath,amssymb,amsthm,physics}
\usepackage{geometry}
\usepackage{hyperref}
\geometry{margin=1in}
\title{Possible Connections Between Nonrelativistic Quantum Mechanics and Einstein--aether Theory}
\author{}
\date{}
\begin{document}
\maketitle
\begin{abstract}
Einstein--aether (EA) theory is a generally covariant modification of
general relativity containing a dynamical unit timelike vector field
(the aether) that defines a locally preferred frame while preserving
diffeomorphism invariance. Nonrelativistic quantum mechanics (NRQM),
despite its enormous empirical success, continues to raise foundational
questions concerning the origin of the Schr\"odinger equation,
measurement, nonlocality, and the role of time.
This article reviews possible conceptual and mathematical links between
NRQM and Einstein--aether theory. Particular attention is given to:
(i) preferred-frame interpretations of quantum theory,
(ii) stochastic and hydrodynamic formulations,
(iii) emergent quantum mechanics,
(iv) Lorentz-violating ultraviolet completions,
and (v) the possibility that the aether field provides a physical
substrate for quantum phenomena.
Although no accepted derivation of quantum mechanics from Einstein--aether
theory exists, several intriguing correspondences suggest directions for
future research.
\end{abstract}
\section{Introduction}
Two major conceptual difficulties of modern theoretical physics remain:
\begin{enumerate}
\item The incompatibility between quantum mechanics and gravity.
\item The absence of a universally accepted interpretation of quantum
mechanics.
\end{enumerate}
Einstein--aether theory introduces a preferred timelike direction
through a dynamical vector field $u^a$ satisfying
\begin{equation}
u^a u_a = -1.
\end{equation}
The existence of such a field breaks local Lorentz symmetry
spontaneously while preserving general covariance.
The theory therefore occupies an intermediate position between
standard general relativity and theories possessing an explicit
preferred frame.
Historically, preferred-frame concepts have repeatedly appeared in
attempts to understand quantum theory, including:
\begin{itemize}
\item de Broglie-Bohm theory,
\item Nelson stochastic mechanics,
\item stochastic electrodynamics,
\item objective collapse models,
\item superdeterministic approaches,
\item emergent quantum mechanics.
\end{itemize}
This motivates investigation of possible relationships between
Einstein--aether theory and nonrelativistic quantum mechanics.
\section{Einstein--aether Theory}
The action is
\begin{equation}
S=
\frac{1}{16\pi G}
\int d^4x \sqrt{-g}
\left(
R + L_{AE}
\right),
\end{equation}
where
\begin{equation}
L_{AE}
=
-K^{ab}_{\ \ mn}
\nabla_a u^m
\nabla_b u^n
+\lambda(u^a u_a+1).
\end{equation}
The tensor
\begin{equation}
K^{ab}_{\ \ mn}
=
c_1 g^{ab}g_{mn}
+c_2\delta^a_m\delta^b_n
+c_3\delta^a_n\delta^b_m
-c_4 u^a u^b g_{mn}
\end{equation}
contains four dimensionless coupling constants.
A useful object is the induced spatial metric
\begin{equation}
h_{ab}=g_{ab}+u_a u_b .
\end{equation}
In the user's notation one often encounters
\begin{equation}
\tilde g_{ab}
=
g_{ab}-2u_a u_b ,
\end{equation}
which effectively changes the sign along the aether direction.
Such structures naturally separate temporal and spatial degrees of
freedom.
\section{The Preferred Time Problem}
One of the oldest puzzles in quantum theory concerns time.
In NRQM the Schr\"odinger equation reads
\begin{equation}
i\hbar \frac{\partial \psi}{\partial t}
=
\hat H \psi.
\end{equation}
Time is absolute and external.
In contrast, general relativity contains no preferred temporal
parameter. The Hamiltonian constraint
\begin{equation}
H=0
\end{equation}
leads to the famous problem of time in quantum gravity.
Einstein--aether theory partially restores an absolute temporal
structure by introducing a distinguished congruence of timelike
worldlines.
One may define a preferred parameter $\tau$ satisfying
\begin{equation}
u^a \nabla_a \tau =1.
\end{equation}
The resulting foliation resembles the Newtonian time entering NRQM.
Consequently, EA theory naturally bridges
\begin{equation}
\text{GR}
\quad\longrightarrow\quad
\text{absolute time}
\quad\longrightarrow\quad
\text{NRQM}.
\end{equation}
This may be its most immediate conceptual connection to
nonrelativistic quantum theory.
\section{Quantum Equilibrium and Preferred Frames}
Bohmian mechanics introduces trajectories
\begin{equation}
\frac{d\mathbf{x}}{dt}
=
\frac{\nabla S}{m},
\end{equation}
where
\begin{equation}
\psi=R e^{iS/\hbar}.
\end{equation}
For many-particle systems nonlocal influences are instantaneous in a
preferred frame.
A long-standing criticism of relativistic Bohmian theories is the need
for a preferred foliation.
Einstein--aether theory naturally supplies such a foliation.
One may identify Bohmian simultaneity surfaces with hypersurfaces
orthogonal to $u^a$:
\begin{equation}
u_a dx^a =0.
\end{equation}
Thus EA theory provides a covariant realization of the preferred frame
already implicit in Bohmian dynamics.
\section{Connection with Nelson Stochastic Mechanics}
Nelson's stochastic mechanics assumes particle trajectories obey
\begin{equation}
dX_t=b(X_t,t)dt+\sqrt{\frac{\hbar}{m}}\,dW_t.
\end{equation}
The Schr\"odinger equation emerges from diffusion together with a
time-symmetric variational principle.
The unexplained element is the origin of the universal noise field.
An intriguing possibility is that fluctuations of the aether generate
effective stochastic forces.
Schematically,
\begin{equation}
u^a
=
\bar u^a + \delta u^a .
\end{equation}
If
\begin{equation}
\langle \delta u^a \rangle =0
\end{equation}
and
\begin{equation}
\langle
\delta u^a(x)
\delta u^b(y)
\rangle
\neq 0,
\end{equation}
then particles coupled to the aether experience random motion.
Under suitable approximations one may obtain an effective diffusion
equation
\begin{equation}
\partial_t \rho
=
D\nabla^2\rho.
\end{equation}
This suggests a possible microscopic origin for Nelson's diffusion
constant.
No complete derivation currently exists.
\section{Hydrodynamic Quantum Mechanics}
The Madelung transformation yields
\begin{equation}
\psi=\sqrt{\rho}e^{iS/\hbar}.
\end{equation}
The Schr\"odinger equation becomes
\begin{equation}
\partial_t\rho
+
\nabla\cdot
\left(
\rho \frac{\nabla S}{m}
\right)
=
0,
\end{equation}
and
\begin{equation}
\partial_t S
+
\frac{(\nabla S)^2}{2m}
+
V
+
Q
=
0,
\end{equation}
where
\begin{equation}
Q
=
-\frac{\hbar^2}{2m}
\frac{\nabla^2\sqrt{\rho}}
{\sqrt{\rho}}
\end{equation}
is the quantum potential.
If the aether behaves as a fluid-like medium, one may ask whether
$Q$ arises from collective properties of this medium.
This idea resembles:
\begin{itemize}
\item stochastic electrodynamics,
\item pilot-wave hydrodynamics,
\item emergent quantum mechanics,
\item analogue gravity.
\end{itemize}
The challenge is reproducing the exact quantum potential without
introducing it by hand.
\section{Einstein--aether Theory and Lorentz-Violating Quantum Dynamics}
The presence of $u^a$ allows Lorentz-violating operators.
For a scalar field,
\begin{equation}
\mathcal L
=
-\frac12
g^{ab}
\partial_a\phi
\partial_b\phi
-\frac{\alpha}{M^2}
(u^a\partial_a)^2\phi
+\cdots .
\end{equation}
The dispersion relation may become
\begin{equation}
\omega^2
=
c^2k^2
+
\frac{k^4}{M^2}
+\cdots .
\end{equation}
Taking the low-energy limit yields an effective Schr\"odinger equation
\begin{equation}
i\hbar \partial_t\psi
=
-\frac{\hbar^2}{2m}
\nabla^2\psi
+
\frac{\hbar^4}{8m^3M^2}
\nabla^4\psi
+\cdots.
\end{equation}
Such corrections could in principle provide observable deviations from
standard NRQM.
\section{Emergent Quantum Mechanics from the aether}
A speculative but attractive possibility is
\begin{equation}
\text{aether dynamics}
\longrightarrow
\text{statistical coarse graining}
\longrightarrow
\text{Schr\"odinger equation}.
\end{equation}
The proposed mechanism resembles the emergence of thermodynamics from
molecular dynamics.
The key requirements would be:
\begin{enumerate}
\item emergence of complex amplitudes,
\item emergence of Born probabilities,
\item emergence of interference,
\item emergence of entanglement,
\item recovery of Bell correlations.
\end{enumerate}
No current Einstein--aether model satisfies all five conditions.
\section{Relation to Ho\v rava Gravity}
The strongest modern connection may arise indirectly through
Ho\v rava gravity.
Ho\v rava gravity introduces a preferred foliation and anisotropic
scaling
\begin{equation}
t\rightarrow b^z t,
\qquad
x\rightarrow bx.
\end{equation}
Its low-energy limit is closely related to Einstein--aether theory.
Many researchers view Ho\v rava gravity as a candidate ultraviolet
completion of gravity.
Since nonrelativistic scaling is fundamental in Ho\v rava gravity,
there exists a natural conceptual bridge between:
\begin{equation}
\text{nonrelativistic quantum systems}
\leftrightarrow
\text{preferred-foliation gravity}.
\end{equation}
This remains an active area of research.
\section{Major Obstacles}
Several obstacles prevent identification of Einstein--aether theory
with a foundation of quantum mechanics.
\begin{enumerate}
\item No derivation of the Schr\"odinger equation exists.
\item No derivation of the Born rule exists.
\item Bell inequality violations must be reproduced.
\item Relativistic quantum field theory must emerge.
\item Experimental constraints strongly restrict Lorentz violation.
\end{enumerate}
These challenges are substantial.
\section{Conclusions}
Einstein--aether theory and nonrelativistic quantum mechanics share a
common structural feature absent in ordinary general relativity:
the existence of a preferred notion of time.
The most promising connections are:
\begin{enumerate}
\item Preferred foliations for Bohmian mechanics.
\item Physical realization of hidden-variable simultaneity.
\item Possible stochastic origin of Nelson diffusion.
\item Emergent hydrodynamic descriptions.
\item Lorentz-violating quantum corrections.
\item Ho\v rava-inspired ultraviolet completions.
\end{enumerate}
At present there is no accepted derivation of quantum mechanics from
Einstein--aether dynamics. Nevertheless, the theory provides one of
the few mathematically consistent frameworks in which preferred-frame
ideas can be investigated without abandoning general covariance.
Consequently it remains a valuable laboratory for exploring whether
quantum mechanics may ultimately emerge from a deeper spacetime
substrate.
\end{document}
